%! Solving by Square Roots & Completing the Square — Unit 3, Lesson 3.3 — Warm-Up (Answer Key)
\documentclass[10pt]{article}
\usepackage{algebra2-article}
\usepackage{algebra2-key}   % pulls in -boxes; do NOT also load -boxes
\newcommand{\TallMath}[1]{$\displaystyle #1\rule[-1.4em]{0pt}{3.2em}$}

\begin{document}

\pageheader{Unit 3, Lesson 3.3}{Warm-Up --- Answer Key}
\namedateperiod

\vspace{0.10in}

\begin{notesbox}{Getting Ready: Skills We'll Use Today}
\small These four skills power today's two methods. Work quickly.

\vspace{4pt}
\textbf{1. Simplify each radical.} Pull out the largest perfect square. (Simplest radical form.)
\[
  \sqrt{49}=\ans{$7$} \qquad \sqrt{12}=\ans{$2\sqrt{3}$} \qquad
  \sqrt{72}=\ans{$6\sqrt{2}$} \qquad \sqrt{45}=\ans{$3\sqrt{5}$}
\]

\vspace{4pt}
\textbf{2. Square each binomial.} Then look at the constant you get.
\[
  (x+5)^2=\ans{$x^2+10x+25$} \qquad (x-3)^2=\ans{$x^2-6x+9$}
\]
Notice: the constant equals $\big(\tfrac{1}{2}\cdot\text{middle coefficient}\big)^2$. Check it on $(x+5)^2$: half of $10$ is $5$, and $5^2=\ans{$25$}$.

\vspace{4pt}
\textbf{3. Complete the square.} What constant makes each a perfect-square trinomial? Then write it as a square.
\begin{center}
\renewcommand{\arraystretch}{1.5}
\begin{tabular}{|l|c|c|}
\hline
 & \textbf{add this constant} & \textbf{write as a square} \\ \hline
$x^2+8x+\underline{\hphantom{00}}$  & \ans{$16$} & \ans{$(x+4)^2$} \\ \hline
$x^2-10x+\underline{\hphantom{00}}$ & \ans{$25$} & \ans{$(x-5)^2$} \\ \hline
\end{tabular}
\end{center}

\vspace{4pt}
\textbf{4. Two answers.} Solve for \emph{every} value of $x$. Remember: a number has two square roots.
\[
  x^2=16:\ x=\ans{$\pm4$} \qquad x^2=49:\ x=\ans{$\pm7$}
\]
\end{notesbox}

\begin{teachernote}
Every item is a direct prerequisite. Item~1 keeps square-root answers in simplest form ($\sqrt{72}=6\sqrt2$
is the sign-trap). Item~2 is the perfect-square pattern read \emph{forward}; the ``half the middle, then
square'' observation \emph{is} the completing-the-square recipe. Item~3 runs that recipe numerically
(constant $=(b/2)^2$: $16$ and $25$). Item~4 drills the $\pm$ --- the single most common lost point today.
Keep $(x+4)^2$ and $(x-5)^2$ on the board for the notes.
\end{teachernote}

\end{document}
